%% ch01.tex — Chapter 1: The Big Data Revolution: Scale, Velocity, and Variety
%% Part of "Data Engineering" textbook

\chapter{The Big Data Revolution: Scale, Velocity, and Variety}
\label{ch:foundations}

\epigraph{\textit{``The world's most valuable resource is no longer oil, but data.''}}{The Economist, 2017}

\vspace{4pt}

\begin{chapterlearning}
After completing this chapter, you will be able to:
\begin{enumerate}
  \item Explain why traditional relational database systems cannot handle
        modern data workloads at scale, and identify the specific limits
        that are exceeded.
  \item Define the three original dimensions of big data (Volume,
        Velocity, Variety) and the two additional dimensions (Veracity,
        Value) that form the extended 5Vs framework.
  \item Classify any given data source as structured, semi-structured,
        or unstructured, and identify the appropriate storage and query
        tool for each category.
  \item Distinguish between vertical and horizontal scaling strategies,
        and explain why distributed, horizontally scaled systems dominate
        big data engineering.
  \item State the CAP theorem, explain the meaning of each property, and
        identify how real-world systems (HBase, Cassandra, PostgreSQL)
        position themselves in the CAP triangle.
  \item Compare batch and stream processing models and select the
        appropriate model for a given engineering requirement.
\end{enumerate}
\end{chapterlearning}

\bigskip

Every 60 seconds in 2024, users upload 500 hours of video to YouTube,
execute more than 5 million Google searches, send over one million
WhatsApp messages, and publish approximately 6,000 posts on X (formerly
Twitter). These figures are not curiosities; they represent a
fundamental transformation in the scale, speed, and diversity of data
that computing infrastructure is expected to handle. This chapter
establishes the conceptual and technical foundation for the entire
book. We begin by examining what makes data ``big,'' explore how that
bigness is formally characterized, and introduce the core theoretical
constraints---the CAP theorem in particular---that shape every
architectural decision in data engineering.

%%====================================================================
\section{Why Data Engineering Matters}
\label{sec:why-de}
%%====================================================================

In 2018, the International Data Corporation (IDC) projected that the
global \term{datasphere}---the total volume of data generated, captured,
copied, and consumed worldwide---would reach \textbf{175 zettabytes} by
2025, representing a 23-fold increase over the 7.2 ZB that existed in
2018 \parencite{IDC2018}. To appreciate this scale, consider that one
zettabyte equals $10^{21}$ bytes, or roughly 250 billion standard 4\TB\
hard drives---more drives than there are stars visible to the naked eye
from any single point on Earth.

\begin{table}[ht]
\centering
\caption{Data volume units and their relationship to everyday storage}
\label{tab:volume-units}
\renewcommand{\arraystretch}{1.3}
\begin{tabular}{L{2.4cm} R{2.4cm} L{6.0cm}}
\toprule
\textbf{Unit} & \textbf{Bytes} & \textbf{Illustrative equivalent} \\
\midrule
Kilobyte (KB)   & $10^{3}$  & A short email (plain text) \\
Megabyte (MB)   & $10^{6}$  & A 3-minute MP3 audio track \\
Gigabyte (GB)   & $10^{9}$  & A feature-length HD film \\
Terabyte (TB)   & $10^{12}$ & \textasciitilde310,000 photos from a 12 MP camera \\
Petabyte (PB)   & $10^{15}$ & All US academic research libraries combined \\
Exabyte (EB)    & $10^{18}$ & All words ever spoken by humans (estimated) \\
\textbf{Zettabyte (ZB)}  & $\mathbf{10^{21}}$ & \textbf{2025 global datasphere} \\
Yottabyte (YB)  & $10^{24}$ & Theoretical future scale \\
\bottomrule
\end{tabular}
\end{table}

Traditional \term{relational database management systems} (RDBMS)---the
dominant paradigm since Edgar Codd's foundational work in 1970---are
extraordinary tools for managing well-structured, bounded datasets. A
modern relational database running on a well-provisioned server can
comfortably manage tens of terabytes and sustain thousands of
transactions per second with full ACID guarantees. For a bank's
transaction ledger, a hospital's patient registry, or an airline's
reservation system, a carefully tuned relational system remains the
correct choice. The problem emerges when data characteristics exceed
what any single machine can handle: when datasets are measured in
petabytes rather than terabytes; when data arrives not from a single
transactional application but from millions of concurrent event streams;
and when the data itself is not neatly organized into rows and columns
but exists as images, audio recordings, sensor telemetry, and free-form
text. In these conditions, the fundamental assumptions of relational
database design break down.

\term{Data engineering} is the discipline concerned with the design,
construction, and maintenance of systems that collect, store, transform,
and serve large-scale data for downstream consumers---analysts, machine
learning models, business intelligence dashboards, and operational
applications. A data engineer builds the infrastructure that makes raw
data useful. This field draws on distributed systems, software
engineering, database theory, and applied computer science. The tools of
modern data engineering---distributed file systems, parallel batch
processors, in-memory computation engines, stream processing frameworks,
and pipeline orchestrators---are the subjects of this book.

\begin{keyinsight}{The Data Engineering Mandate}
Data engineering is not about \emph{generating} insight---that is the
domain of data science and analytics. Data engineering is about making
data \emph{available, reliable, and efficiently accessible} at any
scale. The data engineer builds the plumbing; analysts and scientists
use the water.
\end{keyinsight}

%%====================================================================
\section{Defining Big Data: The 3Vs Framework}
\label{sec:three-vs}
%%====================================================================

The term ``big data'' gained formal structure in 2001 when Doug Laney,
then a research analyst at the META Group, published a brief but
enormously influential research note identifying three dimensions along
which data management challenges had outpaced conventional tools
\parencite{Laney2001}. He labeled these dimensions \emph{Volume},
\emph{Velocity}, and \emph{Variety}---the \term{3Vs of big data}. Eleven
years later, Gartner (which had acquired META Group) codified these
dimensions into a formal definition:

\begin{defbox}{Big Data (Gartner/Beyer \& Laney, 2012)}
\textit{``Big data is high-volume, high-velocity and/or high-variety
information assets that demand cost-effective, innovative forms of
information processing that enable enhanced insight, decision making,
and process automation.''}\\[4pt]
\hfill --- Beyer \& Laney, Gartner Research, 2012 \parencite{GartnerBigData2012}
\end{defbox}

Each of the three dimensions corresponds to a distinct class of
engineering challenge. The McKinsey Global Institute, in its landmark
2011 study, reinforced this framing: big data refers to datasets
``whose size is beyond the ability of typical database software tools to
capture, store, manage and analyze'' \parencite{McKinsey2011}. Let us
examine each dimension in turn.

%--------------------------------------------------------------------
\subsection{Volume: The Scale Dimension}
\label{subsec:volume}
%--------------------------------------------------------------------

\term{Volume} refers to the sheer quantity of data that is generated,
accumulated, and must be stored and processed. Volume is the most
intuitively obvious characteristic of big data: a dataset is voluminous
when it exceeds the capacity of conventional tools---typically, when a
single machine cannot store and process it within an acceptable elapsed
time. The following production examples illustrate current operational
scales:

\begin{table}[ht]
\centering
\caption{Volume at production scale (representative 2023--2024 estimates)}
\label{tab:volume-examples}
\renewcommand{\arraystretch}{1.35}
\begin{tabular}{L{3.2cm} L{7.4cm}}
\toprule
\textbf{Organisation} & \textbf{Volume characteristic} \\
\midrule
Meta (Facebook)   & 500\TB\ of new data ingested per day across all platforms \\
Google            & $>$100\PB\ of data processed internally per day (estimated) \\
Walmart           & $>$2.5\PB\ of customer transaction data in analytical systems \\
CERN (LHC)        & $\approx$15\PB\ of particle collision data generated per year \\
Netflix           & $>$1\PB\ of user interaction events per day \\
NYSE              & $\approx$1\TB\ of trading records generated on a peak day \\
\bottomrule
\end{tabular}
\end{table}

From an engineering perspective, volume creates several concrete
challenges. First, no single server---even the highest-specification
available in 2024, with up to 64\TB\ of RAM---can hold a petabyte-scale
dataset in memory. Second, sequential reads from a single high-speed
SSD array proceed at approximately 7\GB/s; reading one petabyte at that
speed requires roughly 41 hours. Third, ingesting 500\TB\ per day means
sustained write throughput of approximately 5.8\GB/s, which no single
storage device can sustain reliably. The engineering response is
\emph{horizontal distribution}: partition data into blocks and spread
those blocks across a cluster of hundreds or thousands of commodity
machines---the approach formalized by the Google File System
\parencite{GFS2003} and, in open-source form, by Hadoop's HDFS
(discussed in Chapter~\ref{ch:hdfs}).

%--------------------------------------------------------------------
\subsection{Velocity: The Speed Dimension}
\label{subsec:velocity}
%--------------------------------------------------------------------

\term{Velocity} refers to the speed at which data is generated, arrives
at a system, and must be processed or acted upon. High-velocity data is
not merely large; it arrives \emph{continuously} and often demands
responses on timescales far shorter than batch processing can provide.

Consider a few representative velocities:
\begin{itemize}
  \item The Visa payment network processes approximately 24,000
        transactions per second at peak load.
  \item Global IoT deployments generate over one billion device sensor
        readings per hour across industrial, environmental, and consumer
        applications.
  \item A ride-sharing platform (such as Uber or Lyft) receives a GPS
        ping every three seconds from millions of simultaneously active
        drivers---at 5 million drivers, that is over 1.6 million
        location records per second.
  \item Modern stock exchanges publish quote updates for listed
        securities at rates exceeding one million messages per second
        during high-activity sessions.
\end{itemize}

The engineering challenge of velocity is distinct from that of volume:
a system may need to respond to each event within milliseconds (fraud
detection, payment authorization) or may tolerate latencies of seconds
(real-time dashboards). The key constraint is that high-velocity data
\emph{cannot wait} for a nightly batch process---by the time the batch
runs, the window for action has closed. This distinction drives the
separation between \emph{batch} and \emph{stream} processing
architectures, explored in Section~\ref{sec:batch-stream} and in depth
in Chapter~\ref{ch:pipelines}.

\begin{table}[ht]
\centering
\caption{The data processing spectrum by velocity requirement}
\label{tab:velocity-spectrum}
\renewcommand{\arraystretch}{1.35}
\begin{tabular}{L{2.2cm} L{1.6cm} L{3.2cm} L{3.4cm}}
\toprule
\textbf{Processing mode} & \textbf{Latency} & \textbf{Representative tools} & \textbf{Typical use case} \\
\midrule
Batch          & Hours      & Hadoop MapReduce, Hive & Nightly ETL, monthly billing \\
Micro-batch    & Seconds    & Spark Streaming        & Near-real-time dashboards \\
Stream         & Milliseconds & Apache Flink, Kafka Streams & Fraud detection, alerts \\
Real-time      & Sub-ms     & In-memory + Kafka      & Payment authorization \\
\bottomrule
\end{tabular}
\end{table}

%--------------------------------------------------------------------
\subsection{Variety: The Diversity Dimension}
\label{subsec:variety}
%--------------------------------------------------------------------

\term{Variety} refers to the diversity of data formats, types, schemas,
and sources that must be jointly ingested, stored, and analysed. Modern
organisations do not operate a single data source with a uniform schema;
they accumulate data from dozens of systems, each with its own format
and semantics.

Variety manifests along several axes. \emph{Format variety} arises when
data arrives as CSV files, JSON documents, XML feeds, binary Avro
records, columnar Parquet files, DICOM medical images, and MP4 video
simultaneously---all of which must be stored and queried together.
\emph{Source variety} arises when the same analytical question requires
joining data from relational databases, REST APIs, IoT sensor streams,
social media firehoses, and internal event logs. \emph{Semantic variety}
arises when a field named \texttt{date} means the transaction date in
one system, the record-entry date in another, and the customer birth
date in a third.

The foundational distinction among data types---structured,
semi-structured, and unstructured---is the key conceptual framework for
managing variety, and we explore it in detail in
Section~\ref{sec:data-taxonomy}.

\begin{researchspotlight}{Laney (2001) --- The Origin of the 3Vs}
\textbf{Citation:} Laney, D. (2001). \textit{3D Data Management:
Controlling Data Volume, Velocity and Variety}. META Group Research
Note 949, February 2001. \parencite{Laney2001}

\smallskip\noindent
\textbf{Key insight:} Laney's one-page note, written in the context of
e-commerce data challenges, framed data management as a three-dimensional
problem---not merely a storage problem. He argued that organisations
failing to address all three dimensions would face ``increasing
management headaches.''

\smallskip\noindent
\textbf{Technical contributions:}
\begin{itemize}
  \item Established the first formal vocabulary (\emph{3Vs}) for
        characterising modern data management challenges, unifying what
        had previously been treated as unrelated operational problems.
  \item Distinguished \emph{velocity} as a management dimension
        independent of volume---recognising that \emph{how fast} data
        arrives is as important as \emph{how much} arrives.
  \item Introduced \emph{variety} as the hardest dimension to manage,
        predicting the rise of semi-structured and unstructured data as
        dominant forms.
\end{itemize}

\smallskip\noindent
\textbf{Impact today:} The 3Vs framework is taught at every major
computer science programme and has been cited in thousands of academic
papers and industry reports. It remains the canonical definition of big
data, even as the framework has been extended to 5Vs and beyond. Gartner
officially adopted and extended the definition in 2012.
\end{researchspotlight}

%%====================================================================
\section{The Extended 5Vs Framework}
\label{sec:five-vs}
%%====================================================================

The 3Vs framework is powerful but incomplete. Two additional dimensions
have gained wide acceptance, bringing the framework to five:

\term{Veracity} refers to the trustworthiness and quality of data.
High-veracity data is accurate, consistent, and free of critical errors;
low-veracity data is noisy, inconsistent, or of uncertain provenance.
Veracity challenges arise in practical systems for many reasons: GPS
sensors drift by tens of metres in dense urban environments; optical
character recognition (OCR) on historical documents introduces
transcription errors; social media platforms contain significant volumes
of bot-generated content; and sensor networks may experience intermittent
failures that produce spurious or missing readings. The consequences of
low veracity can be severe---a machine learning model trained on
inaccurate data will make inaccurate predictions, regardless of how
sophisticated the model or how large the training set.

\term{Value} refers to the business or societal utility that can be
extracted from data. It is the ultimate justification for the
considerable engineering investment required to build big data systems.
A dataset may score highly on all four other Vs---petabytes of
high-velocity, highly varied, high-quality data---and still produce
no actionable insight if the right analytical questions are never asked,
or if the results are never operationalised. Value is the bridge between
technical capability and organisational outcomes.

\begin{table}[ht]
\centering
\caption{The 5Vs framework with engineering implications}
\label{tab:five-vs}
\renewcommand{\arraystretch}{1.45}
\begin{tabular}{L{1.8cm} L{3.2cm} L{3.5cm} L{2.6cm}}
\toprule
\textbf{Dimension} & \textbf{Definition} & \textbf{Production example} &
\textbf{Primary engineering challenge} \\
\midrule
Volume    & Scale of data generated and stored & Meta: 500\TB/day ingested &
            Distributed storage (HDFS, S3) \\
Velocity  & Speed of generation and processing & Visa: 24,000 tx/s &
            Stream processing (Kafka, Flink) \\
Variety   & Diversity of formats and sources & CSV + JSON + DICOM + MP4 &
            Schema integration, ETL/ELT \\
Veracity  & Trustworthiness of data quality & GPS drift, OCR errors, bot tweets &
            Data quality pipelines, validation \\
Value     & Business utility extracted from data & Fraud prevented, revenue gained &
            Feature engineering, ML pipelines \\
\bottomrule
\end{tabular}
\end{table}

\begin{researchspotlight}{McKinsey Global Institute (2011) --- Quantifying the Value Dimension}
\textbf{Citation:} Manyika, J.\ et~al.\ (2011). \textit{Big Data: The
Next Frontier for Innovation, Competition, and Productivity}. McKinsey
Global Institute, May 2011. \parencite{McKinsey2011}

\smallskip\noindent
\textbf{Key insight:} The report argued that big data has become ``a new
factor of production,'' alongside capital and labour---that is,
organisations with superior data capabilities will, over time,
systematically outperform those without them, in the same way that
access to capital separates competitive enterprises from struggling ones.

\smallskip\noindent
\textbf{Technical and economic contributions:}
\begin{itemize}
  \item Quantified the value potential across five major sectors:
        healthcare (\$300B+ annually in the US), public sector
        (\$250B+ in Europe), retail ($>$60\% operating margin
        improvement for early adopters), manufacturing, and personal
        location data.
  \item Identified the acute shortage of data-skilled talent as the
        primary constraint on value realisation---the report projected
        a shortfall of 140,000--190,000 professionals with deep
        analytical expertise in the US alone by 2018.
  \item Introduced the concept of ``data as infrastructure,'' arguing
        that the economic benefits of big data investment are broadly
        distributed across the economy, much like roads or electricity
        grids.
\end{itemize}

\smallskip\noindent
\textbf{Impact today:} This report is one of the most cited documents
in the history of data-driven business strategy. It directly influenced
technology investment priorities at major corporations and governments
worldwide, and it validated the creation of dedicated data engineering
and data science functions within large organisations.
\end{researchspotlight}

%%====================================================================
\section{A Taxonomy of Digital Data}
\label{sec:data-taxonomy}
%%====================================================================

Managing the Variety dimension of big data requires a clear
classification of data types. The field has converged on three primary
categories---structured, semi-structured, and unstructured---distinguished
principally by the presence, flexibility, and location of their schema.

%--------------------------------------------------------------------
\subsection{Structured Data}
\label{subsec:structured}
%--------------------------------------------------------------------

\begin{defbox}{Structured Data}
Data organised according to a \textbf{predefined, rigid schema} in which
every record contains the same fields in the same order, with each field
having a specified data type. The schema is defined and enforced before
data is written (\textbf{schema-on-write}). Structured data is typically
stored in relational database management systems and queryable directly
with standard SQL.
\end{defbox}

Structured data represents the most mature and well-understood category.
Relational database theory, developed by Codd (1970) and refined over
five decades, provides powerful tools for storing, querying, and
maintaining the integrity of structured data. A bank's transaction ledger
is a canonical example: every row has an account number, a timestamp, an
amount, a currency code, and a transaction type. There are no optional
fields, no embedded documents, no ambiguous formats. SQL can answer
virtually any analytical question over such data with high efficiency.

The primary limitation of structured data is its rigidity. Adding a new
column to a table with hundreds of millions of rows requires a schema
migration---a costly, time-consuming, and potentially disruptive
operation. More fundamentally, many of the most valuable data types in
modern organisations---customer reviews, support transcripts, medical
images, sensor streams---cannot be naturally represented in rows and
columns without destructive simplification. Despite this, structured data
remains the bedrock of transactional systems, and big data engineering
tools like \tool{Apache Sqoop} exist specifically to migrate structured
data from RDBMS sources into distributed analytical systems (covered in
Chapter~\ref{ch:pipelines}).

In terms of global prevalence, structured data accounts for only
approximately 20\% of all digital data---yet it represents the most
well-governed, analysable slice of that data.

%--------------------------------------------------------------------
\subsection{Semi-Structured Data}
\label{subsec:semi-structured}
%--------------------------------------------------------------------

\begin{defbox}{Semi-Structured Data}
Data with a \textbf{flexible, self-describing schema} embedded within
the data itself. Fields are identified by keys or tags that travel with
the record, not by position in a fixed row. New fields can be added
without modifying a central schema registry. Semi-structured data is
typically stored as files or in NoSQL databases and processed with a
\textbf{schema-on-read} approach: the schema is imposed at query time,
not at write time.
\end{defbox}

Semi-structured data dominates modern web and application architectures.
The two most prevalent formats are \textbf{JSON} (JavaScript Object
Notation) and \textbf{XML} (Extensible Markup Language):

\begin{lstlisting}[style=pseudocode, caption={A JSON event record from a streaming API}]
{
  "event_id":   "e-928471",
  "event_type": "page_view",
  "timestamp":  "2024-03-15T14:22:31Z",
  "user": {
    "id":       "u-441982",
    "tier":     "premium"
  },
  "page":       "/product/laptop-x500",
  "session_ms": 4820,
  "metadata": {
    "user_agent": "Mozilla/5.0 ...",
    "referrer":   "https://search.example.com"
  }
}
\end{lstlisting}

Notice that the JSON record carries its own field names (``event\_id'',
``timestamp'', ``user'', etc.) directly within the data. A consumer
reading this record does not need to consult an external schema
definition to know what each field means. This \emph{self-description}
property makes semi-structured data highly flexible: a downstream
service that only cares about \texttt{event\_type} and
\texttt{user.tier} can simply ignore all other fields. When a new field
(e.g., \texttt{ab\_test\_variant}) must be added, it can be appended to
future records without invalidating existing ones.

The schema-on-read approach enabled by semi-structured formats is
central to the \emph{data lake} architecture: data is ingested in its
native format without transformation, and the schema is applied by the
analytical tool at query time. This defers the cost of schema design
but introduces the risk of schema drift---where the meaning or format
of a field changes silently between producers, creating silent data
quality failures. Managing this risk (through schema registries,
contract testing, and data lineage tracking) is a significant part of
production data engineering practice.

Web server access logs, IoT device telemetry published via MQTT,
REST API responses, email headers, configuration files in YAML or TOML,
and Apache Avro records with embedded schemas all belong to the
semi-structured category. This category accounts for roughly 10\% of
global data by volume.

%--------------------------------------------------------------------
\subsection{Unstructured Data}
\label{subsec:unstructured}
%--------------------------------------------------------------------

\begin{defbox}{Unstructured Data}
Data with \textbf{no predefined schema} and no embedded key-value
structure that can be directly queried. The semantic content is encoded
in the raw signal---pixels, audio samples, token sequences---and cannot
be accessed by a relational or document query engine without substantial
preprocessing (feature extraction, natural language processing, or
computer vision).
\end{defbox}

Unstructured data is the dominant form of digital data, accounting for
an estimated 70--80\% of the global datasphere. It encompasses text
documents (news articles, legal contracts, clinical notes, social media
posts), images (medical scans, satellite photography, product
photographs), audio recordings (call-centre conversations, podcasts,
speech-enabled device logs), video (surveillance footage, user-generated
content, manufacturing inspection recordings), and raw sensor streams
(seismic signals, EEG waveforms, raw GPS trajectories).

The defining challenge of unstructured data is that extracting
information requires \emph{interpretation}---a step that structured and
semi-structured data largely skip. A relational database can compute the
average transaction amount in a single SQL aggregation; answering the
question ``what is the sentiment of customer reviews mentioning our new
product?'' requires tokenising text, building or loading a language
model, running inference, and aggregating the model's output. This
computational chain is far more complex and resource-intensive, which is
why the tools introduced in later chapters---Spark MLlib, distributed
deep learning frameworks, and specialised feature stores---are essential
components of modern data engineering systems.

\begin{table}[ht]
\centering
\caption{Comparison of data type characteristics and engineering implications}
\label{tab:data-type-comparison}
\renewcommand{\arraystretch}{1.45}
\begin{tabular}{L{2.0cm} L{2.6cm} L{2.6cm} L{2.8cm}}
\toprule
\textbf{Dimension}       & \textbf{Structured}     & \textbf{Semi-structured}  & \textbf{Unstructured} \\
\midrule
Schema                   & Predefined, rigid        & Self-describing, flexible & Absent \\
Schema timing            & On write                 & On read                   & N/A \\
Storage                  & RDBMS tables             & Files, NoSQL, APIs        & Object stores, HDFS \\
Query tool               & SQL                      & JSONPath, Hive, Spark SQL & NLP, CV, Spark MLlib \\
Global share             & $\sim$20\%               & $\sim$10\%                & $\sim$70--80\% \\
Schema change cost       & High (migration)         & Low                       & N/A \\
Typical ingestion tool   & Sqoop                    & Kafka, Flume              & Custom pipelines \\
\textbf{Examples}        & Bank ledger, ERP records & REST API JSON, server logs & MRI scan, audio call \\
\bottomrule
\end{tabular}
\end{table}

\begin{caution}{JSON is Semi-Structured, Not Unstructured}
A frequent and consequential error: classifying JSON or XML as
``unstructured'' because it does not look like a database table. JSON
and XML carry an embedded schema in the form of key names and
hierarchical nesting. A query engine can navigate this structure
directly (using JSONPath, XPath, or Hive's schema-on-read). An MRI
scan, a JPEG image, or an MP3 recording carries \emph{no} such
navigable schema---its semantic content is encoded in the binary signal
itself, requiring a domain-specific model (computer vision, speech
recognition) to interpret. The distinction matters architecturally:
semi-structured data can be stored and queried in a data lake with
minimal transformation; unstructured data requires a preprocessing
pipeline before it becomes analytically useful.
\end{caution}

%%====================================================================
\section{The Engineering Response to Scale}
\label{sec:engineering-response}
%%====================================================================

Given the 3Vs, the engineering question is: how do we build systems
capable of storing petabytes, processing millions of events per second,
and integrating dozens of heterogeneous data sources? The answer, which
runs through every chapter of this book, reduces to a small set of
foundational principles.

%--------------------------------------------------------------------
\subsection{Vertical vs.\ Horizontal Scaling}
\label{subsec:scaling}
%--------------------------------------------------------------------

There are two ways to make a computing system more powerful.
\term{Vertical scaling} (scaling up) means adding more resources to a
single machine: a faster processor, more RAM, a larger or faster storage
array. Vertical scaling is simple---it introduces no distributed-systems
complexity---but it has a hard physical ceiling. The most powerful
single-node servers available in 2024 top out at approximately 64\TB\
of RAM and a handful of high-throughput storage controllers. Above this
limit, no amount of money can buy a single machine large enough to hold
the dataset. Vertical scaling is also exponentially expensive: a server
with twice the RAM costs approximately eight times as much.

\term{Horizontal scaling} (scaling out) means adding more machines to a
cluster and distributing the data and computation across all of them.
Each individual machine is a commodity server costing between three and
eight thousand dollars. The power of the cluster grows nearly linearly
with the number of nodes added: doubling the number of nodes
approximately doubles the throughput. Failure of a single node---which
is expected, not exceptional, in a large cluster---does not halt the
system, because data is replicated across multiple nodes.

\begin{table}[ht]
\centering
\caption{Vertical vs.\ horizontal scaling: a systematic comparison}
\label{tab:scaling}
\renewcommand{\arraystretch}{1.4}
\begin{tabular}{L{2.8cm} L{3.3cm} L{3.9cm}}
\toprule
\textbf{Attribute}        & \textbf{Vertical (Scale Up)}  & \textbf{Horizontal (Scale Out)} \\
\midrule
Scalability ceiling       & Hard physical limit ($\sim$64\TB\ RAM) & Near-unlimited (add nodes) \\
Cost scaling              & Superlinear (exponential)    & Near-linear per node \\
Fault model               & Single point of failure      & Tolerates individual node failures \\
Complexity                & Low (no distributed systems) & High (consistency, coordination) \\
Latency for small queries & Low                          & Higher (network overhead) \\
Typical cost example      & Oracle Exadata: $>$\$500K/yr & 1,000 nodes $\times$ \$5K = \$5M (one-time) \\
\textbf{Used by}          & Traditional RDBMS, data warehouses & Hadoop, Spark, Cassandra, Kafka \\
\bottomrule
\end{tabular}
\end{table}

%--------------------------------------------------------------------
\subsection{The Commodity Hardware Insight}
\label{subsec:commodity}
%--------------------------------------------------------------------

The observation that underpins all of distributed data engineering is
deceptively simple: \emph{commodity hardware will fail, and systems must
be designed with this expectation from the start}. In a cluster of 1,000
commodity servers, with each server having a disk failure rate of 1\% per
year, we expect approximately 10 disk failures per year---about one every
five weeks. At 10,000 nodes (the scale of Yahoo!'s first production
Hadoop cluster in 2008), this becomes roughly one disk failure per day.
The solution is not to prevent failures---that is neither economically
nor physically feasible---but to design for them. Every piece of data is
stored in multiple copies on different nodes; if one node fails, the data
remains available from the surviving copies, and the system automatically
creates a new replica on a healthy node. This design principle, first
articulated for distributed file systems by Google's engineers
\parencite{GFS2003}, is the cornerstone of HDFS, Kafka, and Cassandra
alike.

%--------------------------------------------------------------------
\subsection{Moving Computation to Data}
\label{subsec:comp-to-data}
%--------------------------------------------------------------------

The second foundational principle of distributed data processing inverts
the traditional computing model. In a conventional system, data is moved
to the processor: a query is issued, data is fetched from storage into
RAM, and the CPU processes it. This model fails at petabyte scale because
the network bandwidth required to move large amounts of data to a central
processor becomes the bottleneck---the network simply cannot move data
fast enough to keep the processor busy.

The solution, pioneered by Google's MapReduce system \parencite{MapReduce2004}
and implemented in every major big data processing framework, is to
\emph{move computation to the data}: send the processing logic (a
function or query) to the node where the relevant data blocks reside,
execute the computation locally, and collect only the (smaller) results.
This principle---\textbf{data locality}---drastically reduces network
traffic and enables the aggregate CPU capacity of an entire cluster to be
applied in parallel to a large dataset.

\begin{systemdesign}{Data Locality in Practice}
In HDFS (Chapter~\ref{ch:hdfs}), a 128\MB\ data block is stored on three
nodes. When a MapReduce or Spark job must process that block, the
scheduler first attempts to assign the task to one of the three nodes
already holding a replica (task-local scheduling). If those nodes are
busy, it next tries a node on the same physical rack (rack-local
scheduling). Only if no rack-local node is available does it resort to
sending the data over the network. In production clusters, over 90\% of
tasks execute with data locality, meaning almost all processing is local
I/O rather than network transfer.
\end{systemdesign}

%%====================================================================
\section{The CAP Theorem}
\label{sec:cap}
%%====================================================================

Building a distributed system---one that spans multiple machines
connected by a network---introduces a fundamental theoretical constraint
that shapes every storage and database design decision in data
engineering. This constraint is captured by the \term{CAP theorem}.

%--------------------------------------------------------------------
\subsection{Statement and Definitions}
\label{subsec:cap-statement}
%--------------------------------------------------------------------

The CAP theorem was conjectured by Eric Brewer in his keynote address at
the ACM Symposium on Principles of Distributed Computing in 2000
\parencite{Brewer2000} and formally proved by Gilbert and Lynch two years
later \parencite{GilbertLynch2002}. It states:

\begin{defbox}{The CAP Theorem (Brewer, 2000; Gilbert \& Lynch, 2002)}
A distributed data system can simultaneously guarantee \textbf{at most
two} of the following three properties:
\begin{enumerate}
  \item \textbf{Consistency (C):} Every read operation returns the most
        recent successfully written value, or an error. All nodes in the
        system see the same data at the same time.
  \item \textbf{Availability (A):} Every request to a non-failed node
        receives a response---not guaranteed to be the most recent
        write, but a response nonetheless.
  \item \textbf{Partition Tolerance (P):} The system continues to
        operate correctly even if an arbitrary number of messages
        between nodes are lost or delayed (a \emph{network partition}).
\end{enumerate}
\end{defbox}

Understanding what each property means in operational terms is essential
for reasoning about system design. \emph{Consistency}, in the CAP sense,
is not the same as the C in ACID transactions (which refers to
maintaining database invariants); CAP consistency is \emph{linearisability}---
the guarantee that every read sees the most recent write, as if all
operations occurred on a single machine. \emph{Availability} means every
request gets a response; it does not mean the response contains current
data. \emph{Partition tolerance} means the system continues to function
when the network splits into disconnected sub-clusters.

%--------------------------------------------------------------------
\subsection{The Practical Implication: Choosing C vs.\ A}
\label{subsec:cap-practice}
%--------------------------------------------------------------------

The critical practical insight is that network partitions in distributed
systems are not theoretical edge cases---they are \emph{inevitable}. Any
system operating across multiple machines connected by a real network
will, at some point, experience a partition: a switch failure, a
misconfigured firewall rule, a miscabling in a data centre, or a brief
network congestion event that causes messages to be delayed or dropped.
Therefore, in any real distributed system, partition tolerance (P) is
not optional. A system that cannot tolerate partitions is effectively a
single-machine system.

Given that P is always required, the CAP theorem reduces to a binary
choice: when a partition occurs, should the system prioritise
\textbf{consistency} (refuse requests that cannot be answered with
current data, guaranteeing accuracy) or \textbf{availability} (answer
every request with possibly stale data, guaranteeing responsiveness)?
This choice is one of the most important architectural decisions in data
engineering.

\begin{figure}[ht]
\centering
\begin{tikzpicture}[
    node distance=0pt,
    corner/.style={font=\small\bfseries\sffamily, text=DEBlue, align=center},
    edge lbl/.style={font=\scriptsize\sffamily, align=center},
    sys/.style={font=\scriptsize\ttfamily, align=center},
    blob/.style={draw=none, rounded corners=3pt, inner sep=3pt, align=center,
                 font=\scriptsize\sffamily}
  ]
  % Triangle vertices
  \coordinate (C) at (0, 5.2);
  \coordinate (A) at (-4.2, 0);
  \coordinate (P) at (4.2, 0);

  % Triangle
  \draw[DEBlue, very thick] (C) -- (A) -- (P) -- cycle;

  % Vertex labels
  \node[corner] at (0, 5.7)    {Consistency (C)};
  \node[corner] at (-4.9, -0.4) {Availability (A)};
  \node[corner] at (4.9, -0.4)  {Partition Tolerance (P)};

  % CA zone (top left edge)
  \node[blob, fill=DELightBlue] at (-2.8, 3.2) {%
    \textbf{CA} systems\\[2pt]
    PostgreSQL, MySQL\\
    (single-datacenter)
  };

  % CP zone (top right edge)
  \node[blob, fill=DELightGreen] at (2.8, 3.2) {%
    \textbf{CP} systems\\[2pt]
    Apache HBase\\
    ZooKeeper
  };

  % AP zone (bottom)
  \node[blob, fill=DELightAccent] at (0, 0.9) {%
    \textbf{AP} systems\\[2pt]
    Apache Cassandra\\
    Amazon DynamoDB
  };

  % Edge descriptors
  \node[edge lbl, rotate=52, text=DEBlue!60] at (-1.8, 2.4)
        {no partition\\tolerance};
  \node[edge lbl, rotate=-52, text=DEBlue!60] at (1.8, 2.4)
        {sacrifices\\availability};
  \node[edge lbl, text=DEBlue!60] at (0, -0.4)
        {sacrifices consistency};
\end{tikzpicture}
\caption{The CAP triangle: every distributed system chooses two of the three properties.
         Because network partitions are unavoidable, real systems choose between CP and AP.}
\label{fig:cap-triangle}
\end{figure}

\begin{table}[ht]
\centering
\caption{CAP positioning of representative data systems}
\label{tab:cap-systems}
\renewcommand{\arraystretch}{1.4}
\begin{tabular}{L{2.8cm} C{0.8cm} L{5.4cm}}
\toprule
\textbf{System} & \textbf{CAP} & \textbf{Design rationale} \\
\midrule
PostgreSQL, MySQL & CA & Designed for single-node or tightly coupled clusters;
                         cannot tolerate true network partitions \\
Apache HBase      & CP & Chooses consistency: if a partition occurs, HBase may
                         become unavailable rather than serve stale data \\
Apache ZooKeeper  & CP & Coordination service where serving stale configuration
                         data would be catastrophic \\
Apache Cassandra  & AP & Chooses availability: always responds, using
                         tunable eventual consistency (read your own writes) \\
Amazon DynamoDB   & AP & Optimised for always-on e-commerce workloads where
                         availability is paramount \\
\bottomrule
\end{tabular}
\end{table}

\begin{researchspotlight}{Brewer (2000) and Gilbert \& Lynch (2002) --- The CAP Theorem}
\textbf{Citations:}
\begin{itemize}
  \item Brewer, E.A.\ (2000). \textit{Towards Robust Distributed Systems}.
        Keynote, PODC 2000. \parencite{Brewer2000}
  \item Gilbert, S.\ \& Lynch, N.A.\ (2002). Brewer's Conjecture and the
        Feasibility of Consistent, Available, Partition-Tolerant Web
        Services. \textit{ACM SIGACT News}, 33(2):51--59.
        \parencite{GilbertLynch2002}
\end{itemize}

\smallskip\noindent
\textbf{Key insight:} Brewer's conjecture, presented as intuition in
2000, was that three desirable properties of distributed systems could
not all be simultaneously guaranteed. Gilbert and Lynch provided a
formal proof two years later, elevating the observation from engineering
folklore to a mathematically rigorous theorem.

\smallskip\noindent
\textbf{Technical contributions:}
\begin{itemize}
  \item Provided the formal definitions of consistency (linearisability),
        availability, and partition tolerance used in distributed systems
        theory.
  \item Proved that no distributed system can achieve all three
        properties simultaneously using an asynchronous network model.
  \item Established the vocabulary (CA, CP, AP) that data architects use
        to describe and compare distributed systems.
\end{itemize}

\smallskip\noindent
\textbf{Nuance and evolution:} The CAP theorem is often oversimplified.
Eric Brewer himself, in a 2012 retrospective, noted that the choice is
not always a binary C-vs-A decision: systems like DynamoDB offer tunable
consistency levels, allowing different operations to trade consistency
for availability at the operation level. The PACELC model (Patterson
et~al., 2012) extends CAP by noting that even without a partition,
systems must trade latency for consistency---a trade-off CAP does not
capture. Despite these refinements, the CAP theorem remains the most
widely cited result in distributed systems theory.
\end{researchspotlight}

%%====================================================================
\section{Batch vs.\ Stream Processing}
\label{sec:batch-stream}
%%====================================================================

The Velocity dimension of big data---the need to process data that
arrives continuously and at high rates---has driven the development of
two fundamentally different processing paradigms that data engineers must
understand and choose between.

%--------------------------------------------------------------------
\subsection{Batch Processing}
\label{subsec:batch}
%--------------------------------------------------------------------

\term{Batch processing} operates on a fixed, bounded dataset accumulated
over a period of time. The data is first stored in full (on disk, in a
distributed file system, or in a data warehouse), and then a
computation---the \emph{batch job}---is run over the entire stored
dataset. Classic examples include nightly financial reconciliation, end-of-month
billing, weekly machine learning model retraining, and annual census
analysis.

Batch processing has several important properties. First, it achieves
very high \emph{throughput}: because the entire dataset is available
before processing begins, the job can be optimised for sequential I/O,
sorted reads, and massive parallelism. Second, batch jobs are inherently
\emph{reproducible}: re-running the same job on the same input data will
always produce the same output---a property critical for auditing and
debugging. Third, batch jobs can be \emph{fault-tolerant} in a simple
way: if a task fails, it can be restarted from the last checkpoint
without affecting the rest of the job. Fourth, and most importantly for
our purposes, batch processing has \emph{latency measured in hours}: a
nightly ETL job that begins at 02:00 and finishes at 06:00 delivers
results four hours after the data cutoff. For applications that need to
respond to data within seconds or milliseconds, this latency is
unacceptable.

%--------------------------------------------------------------------
\subsection{Stream Processing}
\label{subsec:stream}
%--------------------------------------------------------------------

\term{Stream processing} operates on an unbounded, continuously arriving
sequence of data records. Rather than accumulating data and running a
batch job, a stream processing engine processes each record (or a
small micro-batch of records) as it arrives, maintaining running
state across the stream. The key characteristic is \emph{low latency}:
a stream processor can emit results within milliseconds to seconds of
each event's arrival.

Stream processing is the appropriate model for applications where the
value of a response decays rapidly with time: real-time fraud detection
(a stolen card must be blocked before the second transaction clears, not
after nightly batch analysis); live stock market monitoring (an
algorithmic trading system needs the current quote, not yesterday's
close); operational dashboards (a call-centre manager needs to see
current queue depth, not yesterday's average); and anomaly detection in
industrial systems (a turbine failure signal is valuable in real time,
not eight hours later).

Stream processing introduces engineering challenges that batch processing
avoids. Records may arrive \emph{out of order} (a mobile device's
network connection may be intermittent, causing events to be delayed and
arrive after later events). Events may carry an \emph{event time}
(when something happened) that differs from \emph{processing time}
(when the event arrived at the system). Stateful computations---for
example, counting the number of page views per user in the past
60 minutes---must be maintained across an infinite stream, requiring
careful memory management and state backends. Fault recovery requires
mechanisms like \emph{exactly-once processing} semantics, which add
significant implementation complexity.

\begin{table}[ht]
\centering
\caption{Batch vs.\ stream processing: a systematic comparison}
\label{tab:batch-stream}
\renewcommand{\arraystretch}{1.45}
\begin{tabular}{L{2.6cm} L{3.1cm} L{3.9cm}}
\toprule
\textbf{Attribute}      & \textbf{Batch processing}    & \textbf{Stream processing} \\
\midrule
Data model              & Bounded, stored dataset      & Unbounded, continuously arriving events \\
Latency                 & Minutes to hours             & Milliseconds to seconds \\
Throughput              & Very high                    & Moderate (per-event overhead) \\
Fault recovery          & Re-run from checkpoint       & Exactly-once semantics required \\
State management        & Simple (no persistent state) & Complex (stateful operators, windows) \\
Out-of-order data       & Not a concern                & Central challenge (watermarks) \\
Debugging               & Easy (fixed inputs/outputs)  & Hard (live, non-reproducible) \\
Typical tools           & Hadoop MapReduce, Hive, Spark Batch & Kafka, Apache Flink, Spark Streaming \\
\textbf{Use cases}      & ETL, ML training, billing    & Fraud detection, monitoring, alerts \\
\bottomrule
\end{tabular}
\end{table}

\begin{keyinsight}{Lambda and Kappa Architectures}
Many production systems need both low latency \emph{and} high
accuracy---properties that are difficult to achieve simultaneously. The
\textbf{Lambda architecture} (Nathan Marz, 2011) addresses this by
running both a batch layer (for accurate, comprehensive historical
analysis) and a speed layer (for low-latency approximate results) in
parallel, merging their outputs in a serving layer. The cost is
engineering complexity: every transformation must be implemented twice,
in different frameworks. The \textbf{Kappa architecture} (Jay Kreps,
2014) simplifies this by replacing the batch layer with a replay of the
full event log through the stream processor---using Apache Kafka's
message retention to make historical reprocessing possible through the
streaming path. Both architectures are examined in depth in
Chapter~\ref{ch:pipelines}.
\end{keyinsight}

%%====================================================================
\section{Case Study: Walmart's Retail Analytics Platform}
\label{sec:walmart-case}
%%====================================================================

\begin{casestudy}{Walmart --- Petabyte-Scale Retail Intelligence}

\textbf{Background.}
Walmart is the world's largest retailer by revenue, operating over
10,000 stores across 27 countries and serving approximately 230 million
customers per week. The scale of its data generation is correspondingly
vast: every store transaction, supply chain event, website interaction,
and third-party demand signal contributes to what is one of the largest
private-sector data repositories in existence.

\smallskip
\textbf{Scale.}
Walmart's internal data analytics environment holds over \textbf{2.5
petabytes} of customer transaction data in its core analytics systems,
with an additional multi-petabyte tier of raw, unstructured data in its
data lake. The company processes roughly one million customer
transactions \emph{per hour}---approximately 280 transactions per second
on a typical day, with significant spikes during seasonal events
(Thanksgiving weekend in the United States generates transaction volumes
three to four times the average daily rate).

\smallskip
\textbf{Data variety.}
Walmart's analytical systems must integrate across all three data
categories described in this chapter:
\begin{itemize}
  \item \textbf{Structured:} Point-of-sale transaction records
        (item, quantity, price, timestamp, store ID, payment method)
        from 10,000+ stores, all originally stored in Oracle RDBMS systems.
  \item \textbf{Semi-structured:} Web and mobile clickstream events in
        JSON format (product views, search queries, cart additions,
        checkout abandonment events) from Walmart.com.
  \item \textbf{Unstructured:} Customer product reviews, call-centre
        transcript audio, and supplier-submitted product imagery for
        catalogue enrichment.
\end{itemize}

\smallskip
\textbf{Engineering challenges and solutions.}
The 3Vs framework maps directly onto Walmart's operational realities:

\begin{itemize}
  \item \textbf{Volume:} Transaction data accumulated over decades cannot
        fit in any single RDBMS. Walmart migrated its analytical
        workloads to a Hadoop-based data lake, using HDFS for distributed
        storage and \tool{Apache Hive} for SQL-style querying over PB-scale
        datasets.
  \item \textbf{Velocity:} Supply chain decisions---whether to redistribute
        inventory across distribution centres---must respond to real-time
        demand signals, not yesterday's batch report. Walmart deployed
        Apache Kafka for real-time event streaming and uses stream
        processing for demand sensing.
  \item \textbf{Variety:} Integrating transaction records, clickstream
        events, weather data (a known driver of consumer behaviour), and
        social media sentiment requires a data integration pipeline that
        handles schema mismatches, different temporal granularities, and
        inconsistent entity representations (the same product may have
        different identifiers in the RDBMS and the web catalogue).
\end{itemize}

\smallskip
\textbf{Business outcomes.}
Walmart's investment in big data engineering has enabled demand
forecasting accuracy that reduced food waste by 15--20\% in its grocery
divisions (reducing spoilage by anticipating demand more precisely), and
dynamic pricing algorithms that run across its full product catalogue
multiple times per day. During the COVID-19 pandemic, real-time analytics
over social media and purchase patterns allowed Walmart to anticipate
demand surges for sanitisers, masks, and canned goods several days before
they peaked---enabling pre-positioning of inventory that proved critical
when supply chains became constrained.

\smallskip
\textbf{Lesson.}
Walmart's evolution from single-RDBMS analytics to a multi-petabyte
distributed data platform illustrates the core thesis of this book: the
3Vs are not abstract concepts but concrete engineering constraints that
force specific architectural choices. Every tool introduced in subsequent
chapters exists because organisations like Walmart encountered these
constraints and could not solve them with conventional technology.
\end{casestudy}

%%====================================================================
\section*{Chapter Summary}
\label{sec:ch1-summary}
%%====================================================================

This chapter established the conceptual and theoretical foundations for
the rest of the book. The key points are:

\begin{itemize}
  \item The global datasphere is projected to reach 175 ZB by 2025,
        driven by social media, IoT, commerce, and scientific
        instrumentation. Traditional RDBMS cannot manage data at this
        scale.

  \item \textbf{Data engineering} is the discipline of building
        infrastructure that makes large-scale data available, reliable,
        and efficiently queryable for downstream consumers.

  \item The \textbf{3Vs framework} (Laney, 2001) characterises big data
        along three dimensions: \emph{Volume} (scale), \emph{Velocity}
        (speed), and \emph{Variety} (diversity). The extended \textbf{5Vs}
        adds \emph{Veracity} (data quality) and \emph{Value} (business
        utility).

  \item Digital data falls into three categories: \textbf{structured}
        (predefined schema, schema-on-write, $\sim$20\% of global data),
        \textbf{semi-structured} (self-describing schema, schema-on-read,
        $\sim$10\%), and \textbf{unstructured} (no schema, requires
        ML/NLP/CV to interpret, $\sim$70--80\%).

  \item \textbf{Horizontal scaling} (adding commodity nodes) is the
        dominant strategy for big data infrastructure because it is
        cost-effective, fault-tolerant, and has no hard ceiling.
        Moving \textbf{computation to data} (data locality) is the
        principle that makes distributed computation efficient.

  \item The \textbf{CAP theorem} proves that any distributed system can
        guarantee at most two of: Consistency, Availability, and Partition
        Tolerance. Since partitions are unavoidable, the real design choice
        is CP vs.\ AP.

  \item \textbf{Batch processing} offers high throughput and simplicity
        for bounded datasets but incurs high latency. \textbf{Stream
        processing} offers low latency for continuously arriving data but
        adds complexity. Production systems often combine both paradigms
        (Lambda or Kappa architectures).
\end{itemize}

%%====================================================================
\section*{Exercises}
\label{sec:ch1-exercises}
%%====================================================================

\exercisesection{Short Questions}

\sq State the 3Vs of big data and give one production-scale numerical
example for each dimension.

\sq What is the CAP theorem? List the three properties and explain what
each means in operational terms.

\sq A single high-end server in 2024 has approximately 64\TB\ of RAM.
Is a 100\PB\ dataset ``big data'' by the Volume definition? Justify your
answer quantitatively.

\sq Classify each of the following as structured, semi-structured, or
unstructured. Justify each answer in one sentence.
\begin{enumerate}[label=(\alph*)]
  \item A table of stock prices in a PostgreSQL database.
  \item A REST API response in JSON format from a weather service.
  \item An MP3 recording of a customer service call.
  \item An XML document containing a purchase order.
  \item A satellite image of an agricultural field.
\end{enumerate}

\sq What is ``schema-on-read''? How does it differ from ``schema-on-write,''
and in which type of data is each used?

\sq Explain the concept of ``data locality'' in distributed processing.
Why does moving computation to data (rather than data to computation)
improve performance at petabyte scale?

\sq An organisation's fraud detection system must decide whether to
approve a payment within 300 milliseconds. Should it use batch or stream
processing? Explain why.

\sq Apache Cassandra is described as an AP system. What does this mean
concretely? Give a scenario where choosing Cassandra over a CP system
(like HBase) is the correct engineering decision, and a scenario where
it is not.

\sq What is the difference between \emph{event time} and
\emph{processing time} in stream processing? Why does this distinction
matter for computing aggregations over a stream?

\sq Identify the primary V (Volume, Velocity, or Variety) that dominates
each scenario:
\begin{enumerate}[label=(\alph*)]
  \item A genomics research lab accumulating 200 TB of whole-genome
        sequencing data per month.
  \item A payment network processing 24,000 card swipes per second.
  \item An enterprise integrating data from 40 departments, each using
        a different database vendor and schema.
\end{enumerate}

\sq What is the Veracity dimension of the 5Vs framework? Give two
concrete examples of low-veracity data from different domains.

\sq Compare the cost-scaling behaviour of vertical and horizontal scaling.
Why is vertical scaling said to be ``superlinear'' in cost?

\sq What is the Lambda architecture? What problem does it solve, and what
is its primary drawback?

\sq True or false, with a one-sentence justification: ``A relational
database system can always be made to handle big data workloads by
upgrading to a more powerful server.''

\bigskip
\exercisesection{Analytical Problems}

\ap \textbf{Data Volume Estimation.}
A financial services company ingests stock market quote updates at a rate
of 1.2 million messages per second during trading hours (6.5 hours per
day). Each message is 256 bytes on average, and HDFS stores three
replicas of all data. Calculate:
\begin{enumerate}[label=(\alph*)]
  \item The raw volume of quote data generated per trading day (in GB).
  \item The total HDFS storage consumed per trading day (with 3$\times$ replication).
  \item The annual storage requirement (250 trading days per year), in TB.
  \item If each HDFS DataNode has 48 TB of usable disk capacity, what is
        the minimum number of DataNodes needed to store one year of data?
\end{enumerate}

\ap \textbf{CAP Theorem Analysis.}
For each of the following systems, determine whether its designers should
choose CP or AP, and explain the reasoning in terms of what happens when
a network partition occurs and the system must choose between C and A.
\begin{enumerate}[label=(\alph*)]
  \item An inventory management system for an online retailer. Showing
        a product as ``in stock'' when it is actually sold out results in
        order cancellations and customer dissatisfaction; showing it as
        ``out of stock'' when it is available causes lost sales.
  \item A distributed configuration registry used by microservices to
        discover the addresses of other services. Serving a stale
        configuration might cause a service to connect to a
        decommissioned endpoint.
  \item A social media ``like'' counter. The precise count at any given
        millisecond is not critical; users expect the counter to update
        within a few seconds.
\end{enumerate}

\ap \textbf{Batch vs. Stream Latency Analysis.}
A credit card company runs a fraud detection system. The fraudster's
behaviour pattern (multiple small test transactions followed by a large
transaction) is detectable within a 15-minute time window. The company
currently runs a batch fraud detection job that starts at midnight and
finishes by 05:00. Show, with a timeline analysis, why this batch job
cannot prevent in-session fraud. Estimate the maximum percentage of daily
fraud that could be caught by a streaming system with a 10-second
detection latency, assuming fraud transactions are uniformly distributed
throughout the day.

\ap \textbf{Scaling Economics.}
A company must build a system to store and query a 500\TB\ analytical
dataset. Option A is a vertically scaled Oracle Exadata appliance
(licensed at \$1.2M/year, maximum capacity 500\TB). Option B is a
horizontally scaled HDFS cluster of commodity servers at \$6,000 per node
(12\TB\ usable per node), with HDFS's 3$\times$ replication factor.
\begin{enumerate}[label=(\alph*)]
  \item How many DataNodes does Option B require?
  \item What is the hardware capital cost of Option B?
  \item If the dataset is expected to grow at 30\% per year, how many
        years before Option A requires a hardware replacement (assume
        it cannot be expanded beyond 500\TB), and what would Option B's
        cost be to scale to the equivalent capacity?
  \item What non-monetary factors also influence this architectural choice?
\end{enumerate}

\bigskip
\exercisesection{Design Problems}

\dpr\textbf{Architecture Selection for a Global E-commerce Platform.}
You are the lead data engineer at a mid-sized e-commerce company with
operations in 12 countries. The company processes 85 million orders per
year and must support the following data workloads:
\begin{itemize}
  \item Monthly financial reconciliation reports (all orders for the month).
  \item Real-time fraud scoring for each checkout event (must respond
        within 500 ms).
  \item A product recommendation engine that retrains weekly on
        full order and clickstream history.
  \item Live operational dashboards showing current order throughput
        and error rates by country.
\end{itemize}
For each workload, identify: (a) the dominant V dimension it stresses,
(b) whether batch or stream processing is appropriate, and (c) which CAP
trade-off the underlying storage must make. Then propose a single
integrated architecture (a sketch is sufficient) that serves all four
workloads, and identify the single component that is the highest
engineering risk.

\dpr\textbf{Data Classification and Integration Strategy.}
A national weather service wants to build an analytics platform
integrating the following data sources:
\begin{itemize}
  \item Hourly temperature and humidity readings from 12,000 automated
        weather stations (CSV format, streamed via SFTP).
  \item Real-time radar imagery in TIFF format, updated every 6 minutes.
  \item Historical weather observations from 1950 to present, stored in
        a legacy Oracle database with a proprietary schema.
  \item Citizen-submitted weather reports via a mobile app (JSON, free-text
        description field included).
  \item Satellite cloud-cover imagery (binary NetCDF format).
\end{itemize}
For each source: (a) classify the data type (structured/semi-structured/
unstructured), (b) identify the dominant V dimension, and (c) propose
an ingestion mechanism. Then identify the two most difficult integration
challenges when combining these sources into a unified query layer, and
propose mitigations for each.

\bigskip
\exercisesection{Programming Exercises}

\pe \textbf{Data Type Classifier (Python).}
The following list of data source descriptions must be classified as
\texttt{structured}, \texttt{semi\_structured}, or \texttt{unstructured}.
Implement the \texttt{classify\_source()} function using keyword-based
heuristics. Then extend your solution to print the recommended primary
storage technology for each category.

\begin{lstlisting}[style=pyspark, caption={PE1.1: Data type classifier}]
sources = [
    {"name": "Customer transaction records in PostgreSQL",
     "format": "relational_table", "schema": "predefined"},
    {"name": "HTTP server access log",
     "format": "text",   "schema": "implicit"},
    {"name": "Satellite multispectral imagery (GeoTIFF)",
     "format": "binary", "schema": "none"},
    {"name": "Product catalog (XML)",
     "format": "xml",    "schema": "self_describing"},
    {"name": "Customer service call recording (MP3)",
     "format": "audio",  "schema": "none"},
    {"name": "E-commerce clickstream events (JSON)",
     "format": "json",   "schema": "self_describing"},
    {"name": "Annual census data (CSV with fixed columns)",
     "format": "csv",    "schema": "predefined"},
]

STORAGE_MAP = {
    "structured":       "Relational RDBMS (PostgreSQL, MySQL) or columnar DW",
    "semi_structured":  "Data lake (HDFS/S3) + Apache Hive / Spark SQL",
    "unstructured":     "Object store (S3, Azure Blob) + ML feature pipeline",
}

def classify_source(source):
    """
    Returns 'structured', 'semi_structured', or 'unstructured'.
    Hint: use source['format'] and source['schema'] fields.
    """
    # TODO: implement classification logic
    pass

for s in sources:
    label = classify_source(s)
    print(f"Source  : {s['name']}")
    print(f"Category: {label}")
    print(f"Storage : {STORAGE_MAP.get(label, 'Unknown')}")
    print()
\end{lstlisting}

\pe \textbf{CAP Trade-off Decision Matrix (Python).}
Given a set of system requirements, write a function that recommends the
appropriate CAP trade-off and a supporting rationale. Then simulate what
happens under a network partition for each choice.

\begin{lstlisting}[style=pyspark, caption={PE1.2: CAP trade-off recommender}]
requirements = [
    {"system": "Banking ledger (account balances)",
     "stale_read_acceptable": False,
     "unavailability_acceptable": True,
     "description": "Must never show incorrect balance"},
    {"system": "Social media post counter (likes)",
     "stale_read_acceptable": True,
     "unavailability_acceptable": False,
     "description": "Slight delay in count update is tolerable"},
    {"system": "Distributed lock manager for microservices",
     "stale_read_acceptable": False,
     "unavailability_acceptable": True,
     "description": "Two services must never acquire the same lock"},
    {"system": "Product search index",
     "stale_read_acceptable": True,
     "unavailability_acceptable": False,
     "description": "Returning yesterday's index is acceptable"},
]

def cap_recommendation(req):
    """
    Returns a tuple: (cap_choice, rationale_string)
    where cap_choice is 'CP' or 'AP'.

    Logic:
    - If stale reads are NOT acceptable -> CP
      (system must be consistent even if it becomes unavailable)
    - If stale reads ARE acceptable    -> AP
      (system must stay available even with possibly stale data)
    """
    # TODO: implement decision logic and rationale string
    pass

def simulate_partition(cap_choice, system_name):
    """
    Print what happens to the system during a network partition.
    """
    if cap_choice == "CP":
        print(f"  Partition occurs in {system_name}:")
        print(f"  -> System refuses requests on isolated nodes.")
        print(f"  -> Availability drops; consistency maintained.")
    elif cap_choice == "AP":
        print(f"  Partition occurs in {system_name}:")
        print(f"  -> System continues serving requests on isolated nodes.")
        print(f"  -> Data may be stale; eventual consistency applies.")

for req in requirements:
    choice, rationale = cap_recommendation(req)
    print(f"System : {req['system']}")
    print(f"Choice : {choice}  |  Rationale: {rationale}")
    simulate_partition(choice, req['system'])
    print()
\end{lstlisting}

\pe \textbf{Storage Volume Estimator (Python).}
Implement a function that estimates the raw and replicated HDFS storage
requirements for a streaming data source, and classifies the result as
``conventional'' (manageable with a single RDBMS), ``big data''
(requiring distributed storage), or ``hyperscale'' (requiring cloud
object storage tiers in addition to HDFS).

\begin{lstlisting}[style=pyspark, caption={PE1.3: HDFS storage volume estimator}]
def estimate_hdfs_storage(records_per_second, avg_record_bytes,
                          retention_days, replication_factor=3):
    """
    Parameters
    ----------
    records_per_second  : float  -- sustained ingestion rate
    avg_record_bytes    : int    -- average size of one record
    retention_days      : int    -- how long to retain data
    replication_factor  : int    -- HDFS replication (default 3)

    Returns
    -------
    dict with keys:
      raw_tb            : raw data in terabytes
      replicated_tb     : total HDFS storage in terabytes (with replication)
      replicated_pb     : total HDFS storage in petabytes
      classification    : 'conventional' | 'big_data' | 'hyperscale'
    """
    BYTES_PER_TB = 1e12
    BYTES_PER_PB = 1e15

    # TODO: implement computation
    # Steps:
    # 1. Compute bytes per second = records_per_second * avg_record_bytes
    # 2. Compute bytes per day = bytes_per_second * 86400
    # 3. Compute raw_bytes = bytes_per_day * retention_days
    # 4. Compute replicated_bytes = raw_bytes * replication_factor
    # 5. Convert to TB and PB
    # 6. Classify:
    #    < 50 TB   -> 'conventional'
    #    50 TB - 10 PB -> 'big_data'
    #    > 10 PB   -> 'hyperscale'
    pass

scenarios = [
    {"name": "IoT temperature sensors (500K devices, 1 reading/min)",
     "rps": 500_000 / 60, "bytes": 64, "days": 365},
    {"name": "E-commerce clickstream (peak 200K events/s, 90-day retention)",
     "rps": 200_000,      "bytes": 512, "days": 90},
    {"name": "Payment network (25K tx/s, 7-year regulatory retention)",
     "rps": 25_000,       "bytes": 256, "days": 2555},
    {"name": "Video streaming events (50K concurrent sessions, 1 event/s)",
     "rps": 50_000,       "bytes": 1024, "days": 180},
]

for s in scenarios:
    result = estimate_hdfs_storage(s["rps"], s["bytes"], s["days"])
    if result:
        print(f"Scenario     : {s['name']}")
        print(f"Raw storage  : {result['raw_tb']:.1f} TB")
        print(f"With 3x rep. : {result['replicated_pb']:.3f} PB")
        print(f"Class        : {result['classification']}")
        print()
\end{lstlisting}
